\documentclass[12pt]{article}
\usepackage{graphicx} % Required for inserting images
\usepackage[margin=0.5in]{geometry}
\usepackage{amsmath, amssymb}
\usepackage{mathrsfs}


\usepackage{soul}


% Dr. David Brown Commands
\newcommand{\ds}{\displaystyle}
\newcommand{\R}{\mathbb{R}}
\newcommand{\N}{\mathbb{N}}
\newcommand{\Q}{\mathbb{Q}}
\newcommand{\C}{\mathbb{C}}


% Commands to get script r
\usepackage{calligra}
\DeclareMathAlphabet{\mathcalligra}{T1}{calligra}{m}{n}
\DeclareFontShape{T1}{calligra}{m}{n}{<->s*[2.2]callig15}{}
\newcommand{\scripty}[1]{\ensuremath{\mathcalligra{#1}}}

% Unit commands
\newcommand{\degree}{\text{°}}
\newcommand{\unitSpeed}{\frac{\text{m}}{\text{s}}}
\newcommand{\unitAcceleration}{\frac{\text{m}}{\text{s}^2}}

% Constant commands
\newcommand{\avogadro}{6.022\times10^{23}}

\newcommand{\boltzmannConstK}{1.381\times10^{-23}\frac{\text{J}}{\text{K}}}
\newcommand{\planckConst}{6.626\times10^{-34}\text{J}\cdot\text{s}}

\newcommand{\atmP}{1.01325\times10^5\,\text{Pa}}

% Commands to easily write partial derivatives
\newcommand{\partDeriv}[2]{\frac{\partial #1}{\partial #2}}
\newcommand{\doublePartDeriv}[2]{\frac{\partial^2 #1}{\partial^2 #2}}


\title{Problem 1.1}
\author{Gabriel Decker}


\begin{document}



\maketitle
\begin{flushleft}

We're given a definition of the Fahrenheit temperature scale in terms of the Celsius scale.
It's given by two points: the freezing point and the boiling point of water at sea level.
In Celsius, the freezing point of water is 0\degree C.
The freezing point of water in the Fahrenheit temperature scale is defined as 32\degree F.
Based on this, we know that at least, the temperature if Fahrenheit, $T_F$ is defined as $T_F=T_C+32\degree$.
However, we're also given that in Celsius, the boiling point of water is 100\degree C, whereas in Fahrenheit, it's 212\degree F.
Clearly, we can't just do $T_F=T_C+112\degree$, since that contradicts our previous finding.
Therefore, we can introduce a factor that multiplies with $T_C$.
So, use the formula $T_F=aT_C+32\degree$.
This doesn't contradict the freezing point definition since $T_F=a\cdot0+32\degree=32\degree$.
So, we can use the definition of $T_F=212\degree$ and $T_C=100\degree$ at the boiling point to determine $a$.\\~\\

$$T_F=aT_C+32\degree$$
$$\implies 212\degree=(a\cdot100\degree)+32\degree$$
$$\implies a\cdot100\degree=180\degree$$
$$\implies a=\frac{180\degree}{100\degree}$$
$$\implies a=\frac{9}{5}$$\\~\\

Therefore, with $a$ determined, the formula for degrees Celsius to Fahrenheit is:
$$T_F=\frac{9}{5}T_C+32\degree$$\\~\\

The inverse of this (Fahrenheit to Celsius) can be found by solving for $T_C$ in the equation we found.
$$T_F=\frac{9}{5}T_C+32\degree$$
$$\implies \frac{9}{5}T_C=T_F-32\degree$$
$$\implies T_C=(T_F-32\degree)\cdot\frac{5}{9}$$\\~\\


Absolute zero in Celsius is defined as $-273.15\degree$ C. We can find it in Fahrenheit by plugging it into the first formula we found.
$$T_F=\frac{9}{5}T_C+32\degree$$
$$\implies T_F=\frac{9}{5}(-273.15\degree)+32\degree$$
$$\implies T_F=-459.67\degree$$





\end{flushleft}

\end{document}
